\section{Inline Visual Route: Proof and Formalization}

These figures are placed in the main argument as visual proof-of-reading aids: each one separates objects, direction, quantitative or formal anchor, and evidence route.

\subsection{Axiom to theorem}

The diagram below is generated from the r012 figure registry, so its objects, arrows, formula anchor, and evidence link are checked before the release audit can pass.

\begin{figure}[p]
\centering
\resizebox{0.96\textwidth}{!}{%
\begin{tikzpicture}[x=1cm,y=1cm,>=Latex,every node/.style={font=\small}]
  \definecolor{ocBlue}{HTML}{173B57}
\definecolor{ocGold}{HTML}{A77D2A}
\definecolor{ocGreen}{HTML}{4B7F52}
\definecolor{ocRed}{HTML}{8E3B32}
\definecolor{ocGray}{HTML}{ECEFF1}
  \draw[rounded corners=10pt,fill=ocGray!35,draw=ocBlue,line width=0.9pt] (-6.8,-3.4) rectangle (6.8,3.4);
  \node[font=\bfseries\large,ocBlue] at (0,3.0) {Axiom to theorem};
  \node[draw,rounded corners=5pt,fill=blue!7,text width=0.24\textwidth,align=center,minimum height=1.05cm] (a) at (-4.35,1.0) {Axiom set};
  \node[draw,rounded corners=5pt,fill=green!8,text width=0.24\textwidth,align=center,minimum height=1.05cm] (b) at (4.35,1.0) {Theorem route};
  \draw[->,line width=1.0pt,ocGreen] (a.east) -- node[above,fill=ocGray!35,inner sep=2pt,align=center] {derivation} (b.west);
  \node[draw,rounded corners=5pt,fill=orange!10,text width=0.28\textwidth,align=center,minimum height=0.9cm] (r) at (-4.35,-1.45) {proof status};
  \node[draw,rounded corners=5pt,fill=white,text width=0.28\textwidth,align=center,minimum height=0.9cm] (m) at (0,-1.45) {proof status};
  \node[draw,rounded corners=5pt,fill=red!6,text width=0.28\textwidth,align=center,minimum height=0.9cm] (f) at (4.35,-1.45) {formal anchor: \( \Gamma\vdash T \)};
  \draw[->,dashed,ocGold] (r.north) -- (a.south);
  \draw[->,dashed,ocGold] (m.north) -- ($(a)!0.5!(b)$);
  \draw[->,dashed,ocGold] (f.north) -- (b.south);
\end{tikzpicture}}
\caption{Axiom to theorem. The figure demonstrates the reader-facing route from Axiom set to Theorem route; the relevant variable or number is proof status, the formula anchor is \textbackslash\{\}Gamma\textbackslash\{\}vdash T, and the evidence link is the corresponding proof, table, or replay section in the release corpus.}
\label{fig:r012-inline-28b_oc133_inline_figures_proof_route-1}
\end{figure}

\subsection{Lean subset}

The diagram below is generated from the r012 figure registry, so its objects, arrows, formula anchor, and evidence link are checked before the release audit can pass.

\begin{figure}[p]
\centering
\resizebox{0.96\textwidth}{!}{%
\begin{tikzpicture}[x=1cm,y=1cm,>=Latex,every node/.style={font=\small}]
  \definecolor{ocBlue}{HTML}{173B57}
\definecolor{ocGold}{HTML}{A77D2A}
\definecolor{ocGreen}{HTML}{4B7F52}
\definecolor{ocRed}{HTML}{8E3B32}
\definecolor{ocGray}{HTML}{ECEFF1}
  \draw[rounded corners=10pt,fill=ocGray!35,draw=ocBlue,line width=0.9pt] (-6.8,-3.4) rectangle (6.8,3.4);
  \node[font=\bfseries\large,ocBlue] at (0,3.0) {Lean subset};
  \node[draw,rounded corners=5pt,fill=blue!7,text width=0.24\textwidth,align=center,minimum height=1.05cm] (a) at (-4.35,1.0) {Human theorem};
  \node[draw,rounded corners=5pt,fill=green!8,text width=0.24\textwidth,align=center,minimum height=1.05cm] (b) at (4.35,1.0) {Mechanized subset};
  \draw[->,line width=1.0pt,ocGreen] (a.east) -- node[above,fill=ocGray!35,inner sep=2pt,align=center] {formal subset} (b.west);
  \node[draw,rounded corners=5pt,fill=orange!10,text width=0.28\textwidth,align=center,minimum height=0.9cm] (r) at (-4.35,-1.45) {checked lemmas};
  \node[draw,rounded corners=5pt,fill=white,text width=0.28\textwidth,align=center,minimum height=0.9cm] (m) at (0,-1.45) {checked lemmas};
  \node[draw,rounded corners=5pt,fill=red!6,text width=0.28\textwidth,align=center,minimum height=0.9cm] (f) at (4.35,-1.45) {formal anchor: \( L\subseteq T \)};
  \draw[->,dashed,ocGold] (r.north) -- (a.south);
  \draw[->,dashed,ocGold] (m.north) -- ($(a)!0.5!(b)$);
  \draw[->,dashed,ocGold] (f.north) -- (b.south);
\end{tikzpicture}}
\caption{Lean subset. The figure demonstrates the reader-facing route from Human theorem to Mechanized subset; the relevant variable or number is checked lemmas, the formula anchor is L\textbackslash\{\}subseteq T, and the evidence link is the corresponding proof, table, or replay section in the release corpus.}
\label{fig:r012-inline-28b_oc133_inline_figures_proof_route-2}
\end{figure}

\subsection{Finite semantics}

The diagram below is generated from the r012 figure registry, so its objects, arrows, formula anchor, and evidence link are checked before the release audit can pass.

\begin{figure}[p]
\centering
\resizebox{0.96\textwidth}{!}{%
\begin{tikzpicture}[x=1cm,y=1cm,>=Latex,every node/.style={font=\small}]
  \definecolor{ocBlue}{HTML}{173B57}
\definecolor{ocGold}{HTML}{A77D2A}
\definecolor{ocGreen}{HTML}{4B7F52}
\definecolor{ocRed}{HTML}{8E3B32}
\definecolor{ocGray}{HTML}{ECEFF1}
  \draw[rounded corners=10pt,fill=ocGray!35,draw=ocBlue,line width=0.9pt] (-6.8,-3.4) rectangle (6.8,3.4);
  \node[font=\bfseries\large,ocBlue] at (0,3.0) {Finite semantics};
  \node[draw,rounded corners=5pt,fill=blue!7,text width=0.24\textwidth,align=center,minimum height=1.05cm] (a) at (-4.35,1.0) {Model instance};
  \node[draw,rounded corners=5pt,fill=green!8,text width=0.24\textwidth,align=center,minimum height=1.05cm] (b) at (4.35,1.0) {Finite witness};
  \draw[->,line width=1.0pt,ocGreen] (a.east) -- node[above,fill=ocGray!35,inner sep=2pt,align=center] {satisfaction} (b.west);
  \node[draw,rounded corners=5pt,fill=orange!10,text width=0.28\textwidth,align=center,minimum height=0.9cm] (r) at (-4.35,-1.45) {model count};
  \node[draw,rounded corners=5pt,fill=white,text width=0.28\textwidth,align=center,minimum height=0.9cm] (m) at (0,-1.45) {model count};
  \node[draw,rounded corners=5pt,fill=red!6,text width=0.28\textwidth,align=center,minimum height=0.9cm] (f) at (4.35,-1.45) {formal anchor: \( M\models\varphi \)};
  \draw[->,dashed,ocGold] (r.north) -- (a.south);
  \draw[->,dashed,ocGold] (m.north) -- ($(a)!0.5!(b)$);
  \draw[->,dashed,ocGold] (f.north) -- (b.south);
\end{tikzpicture}}
\caption{Finite semantics. The figure demonstrates the reader-facing route from Model instance to Finite witness; the relevant variable or number is model count, the formula anchor is M\textbackslash\{\}models\textbackslash\{\}varphi, and the evidence link is the corresponding proof, table, or replay section in the release corpus.}
\label{fig:r012-inline-28b_oc133_inline_figures_proof_route-3}
\end{figure}

\subsection{Negative control}

The diagram below is generated from the r012 figure registry, so its objects, arrows, formula anchor, and evidence link are checked before the release audit can pass.

\begin{figure}[p]
\centering
\resizebox{0.96\textwidth}{!}{%
\begin{tikzpicture}[x=1cm,y=1cm,>=Latex,every node/.style={font=\small}]
  \definecolor{ocBlue}{HTML}{173B57}
\definecolor{ocGold}{HTML}{A77D2A}
\definecolor{ocGreen}{HTML}{4B7F52}
\definecolor{ocRed}{HTML}{8E3B32}
\definecolor{ocGray}{HTML}{ECEFF1}
  \draw[rounded corners=10pt,fill=ocGray!35,draw=ocBlue,line width=0.9pt] (-6.8,-3.4) rectangle (6.8,3.4);
  \node[font=\bfseries\large,ocBlue] at (0,3.0) {Negative control};
  \node[draw,rounded corners=5pt,fill=blue!7,text width=0.24\textwidth,align=center,minimum height=1.05cm] (a) at (-4.35,1.0) {Claim route};
  \node[draw,rounded corners=5pt,fill=green!8,text width=0.24\textwidth,align=center,minimum height=1.05cm] (b) at (4.35,1.0) {Permuted label};
  \draw[->,line width=1.0pt,ocGreen] (a.east) -- node[above,fill=ocGray!35,inner sep=2pt,align=center] {control contrast} (b.west);
  \node[draw,rounded corners=5pt,fill=orange!10,text width=0.28\textwidth,align=center,minimum height=0.9cm] (r) at (-4.35,-1.45) {expected failure};
  \node[draw,rounded corners=5pt,fill=white,text width=0.28\textwidth,align=center,minimum height=0.9cm] (m) at (0,-1.45) {expected failure};
  \node[draw,rounded corners=5pt,fill=red!6,text width=0.28\textwidth,align=center,minimum height=0.9cm] (f) at (4.35,-1.45) {formal anchor: \( control(T)=0 \)};
  \draw[->,dashed,ocGold] (r.north) -- (a.south);
  \draw[->,dashed,ocGold] (m.north) -- ($(a)!0.5!(b)$);
  \draw[->,dashed,ocGold] (f.north) -- (b.south);
\end{tikzpicture}}
\caption{Negative control. The figure demonstrates the reader-facing route from Claim route to Permuted label; the relevant variable or number is expected failure, the formula anchor is control(T)=0, and the evidence link is the corresponding proof, table, or replay section in the release corpus.}
\label{fig:r012-inline-28b_oc133_inline_figures_proof_route-4}
\end{figure}

\subsection{Assumption register}

The diagram below is generated from the r012 figure registry, so its objects, arrows, formula anchor, and evidence link are checked before the release audit can pass.

\begin{figure}[p]
\centering
\resizebox{0.96\textwidth}{!}{%
\begin{tikzpicture}[x=1cm,y=1cm,>=Latex,every node/.style={font=\small}]
  \definecolor{ocBlue}{HTML}{173B57}
\definecolor{ocGold}{HTML}{A77D2A}
\definecolor{ocGreen}{HTML}{4B7F52}
\definecolor{ocRed}{HTML}{8E3B32}
\definecolor{ocGray}{HTML}{ECEFF1}
  \draw[rounded corners=10pt,fill=ocGray!35,draw=ocBlue,line width=0.9pt] (-6.8,-3.4) rectangle (6.8,3.4);
  \node[font=\bfseries\large,ocBlue] at (0,3.0) {Assumption register};
  \node[draw,rounded corners=5pt,fill=blue!7,text width=0.24\textwidth,align=center,minimum height=1.05cm] (a) at (-4.35,1.0) {Assumption};
  \node[draw,rounded corners=5pt,fill=green!8,text width=0.24\textwidth,align=center,minimum height=1.05cm] (b) at (4.35,1.0) {Proof dependency};
  \draw[->,line width=1.0pt,ocGreen] (a.east) -- node[above,fill=ocGray!35,inner sep=2pt,align=center] {dependency edge} (b.west);
  \node[draw,rounded corners=5pt,fill=orange!10,text width=0.28\textwidth,align=center,minimum height=0.9cm] (r) at (-4.35,-1.45) {open assumption total};
  \node[draw,rounded corners=5pt,fill=white,text width=0.28\textwidth,align=center,minimum height=0.9cm] (m) at (0,-1.45) {open assumption total};
  \node[draw,rounded corners=5pt,fill=red!6,text width=0.28\textwidth,align=center,minimum height=0.9cm] (f) at (4.35,-1.45) {formal anchor: \( A_i\Rightarrow T_j \)};
  \draw[->,dashed,ocGold] (r.north) -- (a.south);
  \draw[->,dashed,ocGold] (m.north) -- ($(a)!0.5!(b)$);
  \draw[->,dashed,ocGold] (f.north) -- (b.south);
\end{tikzpicture}}
\caption{Assumption register. The figure demonstrates the reader-facing route from Assumption to Proof dependency; the relevant variable or number is open assumption total, the formula anchor is A\_i\textbackslash\{\}Rightarrow T\_j, and the evidence link is the corresponding proof, table, or replay section in the release corpus.}
\label{fig:r012-inline-28b_oc133_inline_figures_proof_route-5}
\end{figure}

\subsection{Proof closure}

The diagram below is generated from the r012 figure registry, so its objects, arrows, formula anchor, and evidence link are checked before the release audit can pass.

\begin{figure}[p]
\centering
\resizebox{0.96\textwidth}{!}{%
\begin{tikzpicture}[x=1cm,y=1cm,>=Latex,every node/.style={font=\small}]
  \definecolor{ocBlue}{HTML}{173B57}
\definecolor{ocGold}{HTML}{A77D2A}
\definecolor{ocGreen}{HTML}{4B7F52}
\definecolor{ocRed}{HTML}{8E3B32}
\definecolor{ocGray}{HTML}{ECEFF1}
  \draw[rounded corners=10pt,fill=ocGray!35,draw=ocBlue,line width=0.9pt] (-6.8,-3.4) rectangle (6.8,3.4);
  \node[font=\bfseries\large,ocBlue] at (0,3.0) {Proof closure};
  \node[draw,rounded corners=5pt,fill=blue!7,text width=0.24\textwidth,align=center,minimum height=1.05cm] (a) at (-4.35,1.0) {Local proof};
  \node[draw,rounded corners=5pt,fill=green!8,text width=0.24\textwidth,align=center,minimum height=1.05cm] (b) at (4.35,1.0) {Release claim};
  \draw[->,line width=1.0pt,ocGreen] (a.east) -- node[above,fill=ocGray!35,inner sep=2pt,align=center] {promotion boundary} (b.west);
  \node[draw,rounded corners=5pt,fill=orange!10,text width=0.28\textwidth,align=center,minimum height=0.9cm] (r) at (-4.35,-1.45) {closure status};
  \node[draw,rounded corners=5pt,fill=white,text width=0.28\textwidth,align=center,minimum height=0.9cm] (m) at (0,-1.45) {closure status};
  \node[draw,rounded corners=5pt,fill=red!6,text width=0.28\textwidth,align=center,minimum height=0.9cm] (f) at (4.35,-1.45) {formal anchor: \( status\in\{open,closed\} \)};
  \draw[->,dashed,ocGold] (r.north) -- (a.south);
  \draw[->,dashed,ocGold] (m.north) -- ($(a)!0.5!(b)$);
  \draw[->,dashed,ocGold] (f.north) -- (b.south);
\end{tikzpicture}}
\caption{Proof closure. The figure demonstrates the reader-facing route from Local proof to Release claim; the relevant variable or number is closure status, the formula anchor is status\textbackslash\{\}in\textbackslash\{\}\{open,closed\textbackslash\{\}\}, and the evidence link is the corresponding proof, table, or replay section in the release corpus.}
\label{fig:r012-inline-28b_oc133_inline_figures_proof_route-6}
\end{figure}
